% WWW table style: three rules, bold headers, reference green highlight.
\section{Detailed Theoretical Analysis}
\label{app:theory}

The proof is organized around three reductions: orthogonal routing to
independent principal planes; spectral perturbation to a separated Sylvester
equation; and generalization to an exact Gaussian KL followed by binary-KL
inversion. Their classical ingredients and their exact roles are explicit.
The parameter matrix is unconstrained, and a rank-$r$ optimizer state is not
a rank-$r$ parameter constraint. The table distinguishes sharp atomic
inequalities from the sufficient adaptive-state budgets used to combine them.

\begin{table}[htbp]
\centering\small
\caption{\textbf{Proof structure and sharpness.} Exactness concerns the stated
model or atomic inequality, not universal optimality of the complete optimizer
or of every composed finite-sample bound.}
\label{tab:proof-sharpness}
\setlength{\tabcolsep}{4pt}\renewcommand{\arraystretch}{1.18}
\reviewfit{\begin{tabular}{@{}p{.25\linewidth}p{.31\linewidth}p{.37\linewidth}@{}}
\toprule {\bfseries Reduction / classical tool} & {\bfseries Retained quantity} & {\bfseries Equality or scope}\\
\midrule
CS + Poincar\'e separation  &  QR pullback $L_R=\lambda_1-\lambda_d$  &  Two extremal eigenvectors; dimension-free and attained.\\
Ky Fan + square completion  &  Moving-excess coefficient $(1+a^{-1})/(4\Delta)$  &  Optimal to second order in a two-plane perturbation.\\
Principal-angle Rayleigh quotient  &  PL factor $2\Delta\cos^2R$  &  Rank one at angle $R$ attains it.\\
Householder involution  &  $\|D^2\mathcal C_\epsilon[E,E]\|_{\rm op}\le2\|E\|_F^2/\epsilon$  &  Attained at zero with rank-one $E$.\\
Davis--Kahan / Rosenblum  &  Off-diagonal residual divided by the mixed gap  &  Unit residual constant is exact for one Sylvester entry.\\
Descent + Minkowski  &  Quadratic-root stationarity  &  The reference quadratic attains the zero-defect noise coefficient.\\
Yule--Walker / Polyak  &  Exact modal stationary covariance  &  Additive tangent oracle; no multiplicative-noise identification.\\
Gaussian entropy identity  &  Exact prior--posterior angle cost  &  Identity for every positive two-variance covariance.\\
Seeger--Maurer  &  Exact $\Xi_n$ and binary-KL inverse  &  Bernoulli moment is exact; final confidence bound need not be optimal.\\
\bottomrule
\end{tabular}}
\end{table}
The actual Adam defect and the coupled state Lipschitz constants remain
sufficient upper bounds; sharpness is not claimed for those worst-case
budgets. A positive-gap tracking regime is also kept separate from a
stationarity limit in which that gap may vanish.

\subsection{Notation, filtration, and the population being tracked}
\label{app:population}

Write $\langle A,B\rangle=\operatorname{tr}(A^\top B)$ and let $\|\cdot\|$
denote the Frobenius norm unless the operator norm is explicitly indicated.
The pre-sampling filtration $\mathcal F_t$ contains the data set, parameters,
bases, and all optimizer states before drawing the current mini-batch.
Thus $W_t,U_t,M_{t-1},V_{t-1},H_{t-1}$ are measurable with respect to
$\mathcal F_t$. The sampled gradient $G_t$ generally is not. We use two
covariances:
\begin{align}
 S_t&=\mathbb E[G_t^\top G_t\mid\mathcal F_t],\label{eq:raw-cov}\\
 C_t&=\mathbb E\left[\frac{G_t^\top G_t}{\|G_t\|^2+\varepsilon_c}
                 \,\middle|\,\mathcal F_t\right].\label{eq:normalized-cov}
\end{align}
The implemented normalized tracker is unbiased for the tangent gradient of
$\Phi_t(U)=-\operatorname{tr}(U^\top C_tU)/2$. Its eigenspace is not in
general the leading eigenspace of $S_t$. Normalization preserves the eigenvectors
of a \emph{single} covariance sample; it need not preserve eigenvectors after
expectation. For example, let $G^\top G$ equal $100e_1e_1^\top$ with
probability $0.1$ and $e_2e_2^\top$ with probability $0.9$. With
$\varepsilon_c=0$, the unnormalized population selects $e_1$, whereas the
normalized population selects $e_2$. All tracking statements refer to $C_t$;
raw residual-energy statements refer to $S_t$. They coincide up to scale only
under an additional structural assumption, such as deterministic gradient norm.

For a fixed covariance $C\succeq0$, put
\begin{equation}
 \lambda_1\ge\cdots\ge\lambda_d,\quad
 \Delta=\lambda_r-\lambda_{r+1}>0,\quad
 \Lambda=\lambda_1-\lambda_d,\quad
 P=UU^\top,\quad P_\star=U_\star U_\star^\top.
\end{equation}
Here $U_\star$ spans the leading eigenspace. Let
$\theta_1,\ldots,\theta_r\in[0,\pi/2]$ be the principal angles and define
\begin{equation}
 R(U)=\|\Theta\|_F,\quad \theta_{\max}(U)=\max_i\theta_i,
 \quad s(U)=\|\sin\Theta\|_F,\quad h_C(U)=\Phi(U)-\Phi(U_\star).
\end{equation}
The angle $\theta_{\max}$ controls transversality to the optimal subspace;
$R$ is geodesic distance. They must not be identified with the distance between
two consecutive optimizer bases. Constants in the analysis are specific to a
layer, rank, and time interval. A positive average captured energy across layers
does not verify a uniform lower bound on their individual eigengaps.

\subsection{Quotient geometry and implemented operations}

The canonical quotient metric on $\mathrm{Gr}(r,d)$ has horizontal lift
\begin{equation}
 \mathcal H_U=\{\xi:U^\top\xi=0\},\qquad
 \Pi_U(A)=(I-P)A,\qquad g_U(\xi,\zeta)=\langle\xi,\zeta\rangle.
\end{equation}
Equivalently, for projector tangent vectors the metric is half the ambient
Frobenius inner product. These normalizations agree under
$\dot P=\xi U^\top+U\xi^\top$ \cite{bendokat2024grassmann}.
This metric has nonnegative sectional curvature bounded by $2$; the upper
bound need not be attained in every low-dimensional case
\cite{zimmermann2024curvature}. The injectivity radius is $\pi/2$.
The cut locus relative to $[U_\star]$ consists of subspaces having at least
one principal angle $\pi/2$. Consequently $\theta_{\max}<\pi/2$ ensures
the usual graph chart exists, while $R<\pi/2$ is a sufficient, stronger
condition. Neither condition implies geodesic convexity of $\Phi$ throughout
that region.

\begin{lemma}[Projection and exact energy accounting]
\label{lem:decomposition}
For every $G$ and orthonormal $U$,
\begin{equation}
 G=GP+G(I-P),\qquad
 \|G\|^2=\|GU\|^2+\|G(I-P)\|^2.
\end{equation}
For a covariance $S=\mathbb E[G^\top G]$, an
$\mathcal F_t$-measurable projector satisfies
\begin{equation}
 \mathbb E_t\|G(I-P)\|^2
 =\operatorname{tr}(S)-\operatorname{tr}(U^\top SU)
 =\tau_r(S)+2h_S(U),\label{eq:tail-exact}
\end{equation}
where $\tau_r(S)=\sum_{j>r}\lambda_j(S)$.
\end{lemma}
\begin{proof}
The two gradient components are orthogonal because $P(I-P)=0$. Also
$\|GP\|^2=\operatorname{tr}(U^\top G^\top GU)=\|GU\|^2$.
Taking conditional expectation yields the first equality in
\eqref{eq:tail-exact}; adding and subtracting the leading $r$ eigenvalues
yields the second. If $P$ is chosen from the \emph{current} sampled gradient,
the conditional-expectation step requires separate justification.
\end{proof}

\begin{lemma}[QR retraction and its quotient error]
\label{lem:qr}
Let $\xi\in\mathcal H_U$ have compact singular value decomposition
$\xi=A\operatorname{diag}(\sigma_i)B^\top$. The projector of
$\mathrm{qf}(U+\xi)$ is
\begin{equation}
 P_{\mathrm{QR}}=(U+\xi)(I+\xi^\top\xi)^{-1}(U+\xi)^\top.
 \label{eq:qr-projector}
\end{equation}
In the principal planes of $\xi$, QR and the exponential map rotate through
$\arctan\sigma_i$ and $\sigma_i$, respectively. If $\|\xi\|\le1/4$,
\begin{equation}
 d_{\mathrm{Gr}}([\mathrm{qf}(U+\xi)],\operatorname{Exp}_{[U]}\xi)
 \le\frac13\Big(\sum_i\sigma_i^6\Big)^{1/2}
 \le\frac13\|\xi\|^3.\label{eq:qr-exp}
\end{equation}
Projection transport $\mathcal T_{U\to V}\xi=(I-VV^\top)\xi$ is
nonexpansive and satisfies
\begin{equation}
 \|\mathcal T_{U\to V}\xi-\xi\|
 \le\|VV^\top-UU^\top\|_{\mathrm{op}}\|\xi\|.
\end{equation}
\end{lemma}
\begin{proof}
Horizontal tangency gives $(U+\xi)^\top(U+\xi)=I+\xi^\top\xi$,
which is positive definite for every $\xi$. The orthogonal projector onto its
column span is therefore \eqref{eq:qr-projector}. The polar representative
$(U+\xi)(I+\xi^\top\xi)^{-1/2}$ has exactly the same span as its QR
representative. In the $i$th principal plane its coefficients are
$(1,\sigma_i)/\sqrt{1+\sigma_i^2}$, giving angle
$\arctan\sigma_i$. The exponential coefficients are
$(\cos\sigma_i,\sin\sigma_i)$. The small-step condition prevents angle
wrapping. Since
$0\le x-\arctan x=\int_0^x t^2/(1+t^2)\,dt\le x^3/3$,
summing squared angular differences proves \eqref{eq:qr-exp}.
Finally, orthogonal projection cannot increase norm and
$(I-VV^\top)\xi-\xi=(UU^\top-VV^\top)\xi$.
\end{proof}

This calculation identifies the retraction error on the quotient, where the
algorithm actually tracks a subspace. The QR basis itself can also contain a
within-subspace rotation. That rotation affects coordinatewise Adam even when
it has no effect on the routing objective. Numerical orthogonality error,
retraction approximation error, and statistical tracking error are three
different quantities; a small value of the first does not establish either of
the others.

\subsection{The routing landscape: exact principal-angle analysis}
\label{app:landscape}

\begin{theorem}[Landscape, local PL inequality, and its boundary]
\label{thm:landscape}
For $C\succeq0$ with gap $\Delta>0$, the unique minimizing subspace of
$\Phi(U)=-\operatorname{tr}(U^\top CU)/2$ is $[U_\star]$.
Its horizontal gradient and Hessian quadratic form are
\begin{align}
 \operatorname{grad}\Phi(U)&=-(I-P)CU,\label{eq:grass-grad}\\
 \langle\xi,\operatorname{Hess}\Phi(U)[\xi]\rangle
 &=\operatorname{tr}(\xi^\top\xi\,U^\top CU)
   -\operatorname{tr}(\xi^\top C\xi).\label{eq:grass-hess}
\end{align}
Globally,
\begin{equation}
 \frac\Delta2s(U)^2\le h_C(U)\le\frac\Lambda2s(U)^2,
 \qquad \frac2\pi R(U)\le s(U)\le R(U).\label{eq:gap-bounds}
\end{equation}
Whenever $\theta_{\max}(U)\le R_0<\pi/2$,
\begin{equation}
 \|\operatorname{grad}\Phi(U)\|^2
 \ge2\Delta\cos^2 R_0\,h_C(U).\label{eq:PL}
\end{equation}
At the minimum the Hessian eigenvalues are
$\lambda_i-\lambda_j$, $i\le r<j$. Every nonoptimal critical subspace
has a tangent direction of curvature at most $-\Delta$.
Quadratic growth in \eqref{eq:gap-bounds} and PL in \eqref{eq:PL} do
\emph{not} assert geodesic strong convexity on the entire normal neighborhood.
\end{theorem}

\begin{proof}
The variational characterization of eigenvalues gives
$\max_{U^\top U=I}\operatorname{tr}(U^\top CU)=\sum_{i\le r}\lambda_i$.
The gap separates this maximizing subspace from its complement, giving
uniqueness on the quotient. Differentiating $\Phi$ in the ambient space gives
$-CU$; horizontal projection proves \eqref{eq:grass-grad}.
For a Grassmann geodesic with $\dot U(0)=\xi$, one can choose a horizontal
lift satisfying $\ddot U(0)=-U\xi^\top\xi$. Differentiating
$-\operatorname{tr}(U(t)^\top CU(t))/2$ twice proves
\eqref{eq:grass-hess}.

Choose principal vectors $x_i\in[U_\star]$ and
$y_i\in[U_\star]^\perp$ so the columns of a representative of $U$ are
\begin{equation}
 u_i=x_i\cos\theta_i+y_i\sin\theta_i.
\end{equation}
Terms with $\theta_i=0$ need no complementary vector; when the complement
has dimension smaller than $r$, only the nonzero-angle pairs are used.
The $x_i$ form an orthonormal basis for $[U_\star]$ and the relevant
$y_i$ are orthonormal. Since this leading subspace is invariant under $C$,
$x_i^\top Cy_j=0$. Put $a_i=x_i^\top Cx_i$,
$b_i=y_i^\top Cy_i$, and $c_i=a_i-b_i$. Then
\begin{equation}
 h_C(U)=\frac12\sum_i c_i\sin^2\theta_i,
 \qquad \Delta\le c_i\le\Lambda.\label{eq:CS-excess}
\end{equation}
This is an exact excess formula in principal coordinates. The coefficients
$c_i$ are Rayleigh quotients, not generally individual eigenvalue
differences. The bounds on $c_i$ prove the first part of
\eqref{eq:gap-bounds}; $2x/\pi\le\sin x\le x$ on $[0,\pi/2]$
proves the second.

For every nonzero angle set
$e_i=-x_i\sin\theta_i+y_i\cos\theta_i$. The matrix tangent directions
which place $e_i$ in column $i$ and zero elsewhere are orthonormal and
horizontal. Direct differentiation gives their gradient components
$c_i\sin\theta_i\cos\theta_i$. Bessel's inequality therefore gives
\begin{align}
 \|\operatorname{grad}\Phi(U)\|^2
 &\ge\sum_i c_i^2\sin^2\theta_i\cos^2\theta_i\\
 &\ge\Delta\cos^2R_0\sum_i c_i\sin^2\theta_i
 =2\Delta\cos^2R_0\,h_C(U).
\end{align}
This proves PL without a curvature-comparison shortcut. It is consistent
with prior analyses of the block Rayleigh quotient
\cite{alimisis2024geodesic}; the angular factor is not claimed as a new
standalone manifold theorem.

At $U_\star$, diagonalize $C$ in its invariant subspaces. A tangent
coordinate moving selected eigenvector $i$ toward unselected eigenvector
$j$ has Hessian eigenvalue $\lambda_i-\lambda_j$. For a critical point,
$(I-P)CU=0$, so its selected subspace is invariant. If it is not leading,
some selected eigenvector has eigenvalue at most $\lambda_{r+1}$ and some
unselected one has eigenvalue at least $\lambda_r$. Rotating the former
toward the latter gives Hessian curvature at most $-\Delta$.
Finally, for $r=1,d=2$,
\begin{equation}
 h(\theta)=\frac\Delta2\sin^2\theta,\qquad
 h''(\theta)=\Delta\cos(2\theta).
 \label{eq:nonconvex-counterexample}
\end{equation}
Thus the Hessian is negative for $\pi/4<\theta<\pi/2$ even though
quadratic growth and the positive PL bound still hold.
\end{proof}

The distinction in the last sentence is essential for tracking and stability.
PL controls objective decrease along one trajectory. Contraction of the
distance between \emph{two} trajectories requires additional hypotheses;
it does not follow from PL alone. Nor does \eqref{eq:PL} prove that the
scalar $\Delta\cos^2R_0$ predicts the ranking of complete neural-network
optimizers. Gradient noise, covariance normalization, parameter-state
conditioning, and computational cost remain independent factors.

\subsection{Coordinate transport and the limits of diagonal Adam}
\label{app:transport}

\begin{lemma}[First moments and diagonal second moments]
\label{lem:moments}
Let $U_+=UQ$ for $Q\in O(r)$, representing a pure basis change. Then
$Z_+=ZQ$ and exact first-moment consistency requires $M_+=MQ$.
For a row-wise full second moment $B_i=\mathbb E[z_i^\top z_i]$,
the exact transformation is $B_{i,+}=Q^\top B_iQ$. If only its diagonal
$v_i$ is stored, the diagonal-covariance approximation is
\begin{equation}
 v_{i,+}=v_i Q^{\odot2},\label{eq:diag-transport}
\end{equation}
and its error is precisely
\begin{equation}
 \operatorname{diag}(Q^\top B_iQ)-v_iQ^{\odot2}
 =\operatorname{diag}\{Q^\top(B_i-\operatorname{Diag}v_i)Q\}.
 \label{eq:offdiag-error}
\end{equation}
The error need not be second order in the basis rotation. Furthermore,
coordinatewise Adam is not invariant under general $Q$, even with an
exactly transformed diagonal.
\end{lemma}
\begin{proof}
$G(UQ)=(GU)Q$ gives the first claim and
$(MQ)(UQ)^\top=MU^\top$ verifies the lifted first moment.
For second moments, expand $(z_iQ)^\top(z_iQ)$ and take expectation.
Its diagonal under the approximation $B_i=\operatorname{Diag}v_i$ is
$\sum_k v_{ik}Q_{kj}^2$, proving \eqref{eq:diag-transport} and
\eqref{eq:offdiag-error}. For $Q=I+\epsilon A+O(\epsilon^2)$,
the omitted off-diagonal covariance can contribute a linear term through
$\operatorname{diag}(B_iA-AB_i)$.
Finally, entrywise division by a diagonal square root does not commute with
orthogonal mixing. For example, the zero-history, negligible-$\varepsilon$
Adam direction for $z=(2,1)$ is $(1,1)$. A generic rotation followed by
coordinatewise normalization and rotation back does not reproduce that
vector. Signed coordinate permutations are a special invariant case.
\end{proof}

For consecutive, genuinely different subspaces we use
$R=\operatorname{Polar}(U_+^\top U)$ and align $M_+=MR^\top$.
This minimizes basis mismatch, but cannot restore the part of a historical
ambient moment orthogonal to the new span. If an ambient moment is $A$,
then
\begin{equation}
 \|AP-AP_+\|\le\|A\|\|P-P_+\|_{\mathrm{op}}
 =\|A\|\sin\theta_{\max}(U,U_+).
\end{equation}
This is a first-order perturbation bound, not an exact transport identity.
The covariance error from omitted correlations must be added separately.

The archived implementation used $(\sqrt v\,Q)^{\odot2}$ rather than
\eqref{eq:diag-transport}. These differ even for diagonal covariance:
with $v=(1,1)$ and a $45$-degree rotation the former gives $(2,0)$ up to
column order, while the correct diagonal is $(1,1)$. New experiments use
\eqref{eq:diag-transport}; historical outcomes retain their original
algorithm label. Resetting $v$ to zero is another distinct algorithm. A
reset is a deterministic state transformation, not a generic additive
statistical error, and its effect on descent must be analyzed through the
resulting parameter direction.
